\scanpage{143}
\begin{center}- 130 -\end{center}
\margin{n\textsuperscript{o} 316}

invariantes à gauche $X_{ij}$ telles que $X_{ij}^{(e)}=D_{ij}^{(e)}$ sont données par
\[
X_{ij}=(1+U_{ii})D_{ij}+\sum_{k\ne i}U_{ki}D_{kj}.
\]
L'algèbre de Lie de $\mathrm{GL}(n,K)$ s'identifie à $\unsure{\mathrm{gl}|\mathfrak{gl}}(n,K)$ en identifiant $X_{ij}$ à élément $E_{ij}$ de la base canonique (utiliser la formule (3) de l'exercice.24).

Si $K$ est de caractéristique $p>0$, on a $X_{ii}^p=X_{ii}$ et $X_{ij}^p=0$ pour $i\ne j$.

\margin{\P}
26) - a) Soient $K$ un corps de caractéristique $\ne2$, $E$ un espace vectoriel de dimension $n$ sur $K$, $\Phi$ une forme bilinéaire symétrique non dégénérée sur $E$. On suppose $E$ muni d'une base orthogonale pour $\Phi$ et on désigne par $R$ la matrice (diagonale) de $\Phi$ par rapport à cette base. On sait (\uline{Alg}. chap.IX, \S\quad, exerc.\quad) que toute matrice (relative à la base considérée) $U$ du groupe orthogonal $G(\Phi)$ telle que $\det(I+U)\ne0$ s'écrit $U=(I-R^{-1}S)^{-1}(I+R^{-1}S)$, où $S=R(U-I)(U+I)^{-1}$ est une matrice alternée telle que $\det(I-R^{-1}S)\ne0$ ; inversement pour toute matrice $S$ satisfaisant à ces conditions, $U=(I-R^{-1}S)^{-1}(I+R^{-1}S)$ est une matrice de $G(\Phi)$ telle que $\det(I+U)\ne0$. Soient $S_{ij}$, $S'_{ij}$, $S''_{ij}$ $3n(n-1)/2$ indéterminées $(1\leq i<j\leq n)$ et soit $L$ un corps contenant l'anneau de séries formelles $K[[S_{ij},S'_{ij},S''_{ij}]]$. Si on désigne par $S$, $S'$, $S''$ les matrices $\sum_{i<j}S_{ij}(E_{ij}-E_{ji})$, $\sum_{i<j}S'_{ij}(E_{ij}-E_{ji})$, $\sum_{i<j}S''_{ij}(E_{ij}-E_{ji})$, on a $\det(I-R^{-1}S)\ne0$, et les analogues pour $S'$ et $S''$ ; si on pose $U=(I-R^{-1}S)^{-1}(I+R^{-1}S)$, $U'=(I-R^{-1}S')^{-1}(I+R^{-1}S')$, on a aussi $\det(I+UU')\ne0$. Posons $F(S,S')=R(UU'-I)(UU'+I)^{-1}=(f_{ij}(S,S'))$ ; les $f_{ij}(S,S')$ appartiennent à l'anneau des séries formelles
\margin{142}
